\section{Mínimo Produto Víavel (MVP)}

Esta seção tem por objetivo informar as alterações, descartes e escolhas feitas no projeto para o MVP, em relação ao sistema da \acs{poc}. Essas mudanças refletem os problemas e inconsistências enfrentadas durante o desenvolvimento da \acs{poc}, e foram motivadas pela necessidade de otimizar a eficiência e o desempenho do sistema, de forma a atender os requisitos do \acs{mvp} da melhor maneira possível.

\subsection{Mudanças na Infraestrutura}
 
A principal atualização feita no projeto em relação à \acs{poc} foi o descarte do Google Gemini, substituído pelo serviço da OpenAI. Inicialmente, o Gemini foi escolhido devido ao seu baixo custo e à integração facilitada com o Firebase, visando manter todos os serviços em um único ecossistema. Além disso, o Gemini oferecia um serviço de treinamento de modelos customizados, baseado em outros modelos de processamento de texto pré-treinados da Google.
 
Um dos primeiros problemas enfrentados com o serviço do Gemini foi justamente a customização dos modelos. Embora tivéssemos sucesso em treinar e utilizar um modelo com textos de editais e retorno de \acs{json}, não conseguimos utilizar o mesmo fora da plataforma da Google, ou seja, não conseguimos integrar o serviço ao nosso código. Além disso, por ser um serviço relativamente recente, a documentação em certas funcionalidades ainda é escassa, assim como o suporte da comunidade.
 
Os problemas que mais nos afetaram foram a inconsistência e os bloqueios ao utilizar um modelo de processamento de texto já treinado. Os textos \acs{json} gerados pela \acs{ia} eram frequentemente quebrados ou não formatados corretamente, impossibilitando seu consumo dentro da plataforma. Além disso, ao enviar o mesmo texto do edital diversas vezes para processamento, o sistema retornava erros de "recitação" ou de conteúdo não autorizado. Esses erros também possuíam uma documentação extremamente limitada, dificultando ainda mais a solução dos problemas.
 
Essas dificuldades nos levaram a substituir o Google Gemini pelo serviço da OpenAI, que oferecia maior estabilidade e consistência na geração dos conteúdos, melhor documentação e suporte mais robusto. A substituição dos serviços foi um processo relativamente simples, e desde o início os modelos da OpenAI se mostraram muito eficientes. Porém, tivemos a mudança nos custos também, anteriormente, pois estávamos trabalhando com os serviços gratuitos do Gemini, e agora temos outros valores e custos consideravelmente maiores para levar em consideração para o projeto no futuro.
 
A mudança do serviço de \acs{ia} utilizado foi a única alteração relevante nesse período. O restante do sistema e outras integrações se mostraram satisfatórios, atendendo às expectativas e requisitos estabelecidos durante a fase de planejamento.

 \begin{figure}[!htb]
 	    \centering
 	    \caption{\label{logo}QRCode do Deploy: \\ https://studyflow-three.vercel.app}
 	    \includegraphics[width=5cm]{img/qrcode-plataforma.png}
 	    \fonte{Os autores.}
\end{figure}
\FloatBarrier